\begin{thebibliography}{9}
\bibitem{Allouche}
J.-P. Allouche, \emph{Sur la conjecture de ``Syracuse--Kakutani--Collatz''},
S\'eminaire de Th\'eorie des Nombres de Bordeaux (1978--1979),
expos\'e 9, 9-01--9-15, 1979.
\url{https://eudml.org/doc/182044}.

\bibitem{Everett}
C. J. Everett, \emph{Iteration of the number-theoretic function
$f(2n)=n$, $f(2n+1)=3n+2$}, Advances in Mathematics
\textbf{25} (1977), no.~1, 42--45.

\bibitem{Korec}
I. Korec, \emph{A density estimate for the $3x+1$ problem},
Mathematica Slovaca \textbf{44} (1994), no.~1, 85--89.
\url{https://dml.cz/handle/10338.dmlcz/133225}.

\bibitem{Tao7}
T. Tao, \emph{Almost all orbits of the Collatz map attain almost bounded
values}, arXiv:1909.03562v7, 16 July 2026. First submitted 8 September 2019.
\url{https://arxiv.org/abs/1909.03562v7}.

\bibitem{TaoJournal}
T. Tao, \emph{Almost all orbits of the Collatz map attain almost bounded
values}, Forum of Mathematics, Pi \textbf{10} (2022), e12, 56 pp.

\bibitem{Terras}
R. Terras, \emph{A stopping time problem on the positive integers},
Acta Arithmetica \textbf{30} (1976), no.~3, 241--252.
\url{https://doi.org/10.4064/aa-30-3-241-252}.
\end{thebibliography}
